%%%%%%%%%%%%%%%%%%%%%%%%%%%%%%%%%%%%%%%%%%%%%%
%% 01_what_is_cta.tex --- 正文层 (BODY LAYER)
%%
%% 这里只有内容。字体、边距、颜色、定理环境长什么样，全在
%% ../../preamble.tex 里；改版式不用碰这个文件。
%%
%% 这是正文片段，没有 \documentclass 也没有 document 环境：
%% ../../build.sh 单独编译它，all_chapters.tex 把它 \input 进合订本。
%%%%%%%%%%%%%%%%%%%%%%%%%%%%%%%%%%%%%%%%%%%%%%

\chapterhead{01}{What Is a CTA Strategy}{01_what_is_cta}


\section{What a CTA Trades}\label{sec:what-a-cta-trades}

\begin{definition}[CTA strategy]\label{def:cta}
A \emph{CTA (Commodity Trading Advisor) strategy} is a rule-based strategy that trades
\textbf{futures} systematically across asset classes---equities, fixed income, commodities,
currencies---targeting \textbf{absolute return}, independent of market direction.
\end{definition}

\begin{note}
The name is a historical artifact. Modern CTAs are not commodity-only; typical markets are
equity-index futures, Treasuries, FX, energy, metals, and agriculture.
\end{note}

\subsection{Index, ETF, or future}\label{sec:index-etf-future}

``S\&P 500'' names three different objects. Only two of them can be traded.

\begin{tabularx}{\linewidth}{|>{\raggedright\arraybackslash}p{0.14\linewidth}|>{\raggedright\arraybackslash}p{0.18\linewidth}|X|>{\raggedright\arraybackslash}p{0.13\linewidth}|}
\hline
\textbf{Layer} & \textbf{Example} & \textbf{What it is} & \textbf{Tradeable} \\ \hline
\textbf{Index}  & SPX, NDX      & A published number computed from its constituents & No  \\ \hline
\textbf{ETF}    & SPY, VOO, QQQ & A fund holding the constituents; its own shares trade & Yes \\ \hline
\textbf{Future} & ES, NQ        & A contract to settle at a future date; holds nothing & Yes \\ \hline
\end{tabularx}

\begin{definition}[Index]\label{def:index}
\[
  I_t = \frac{M_t}{D_t}, \qquad\qquad M_t = \sum_i N_i P_{it}
\]
where $P_{it}$ is the price of constituent $i$, $N_i$ its float-adjusted share count, $M_t$ the
total float capitalisation, and $D_t > 0$ the \emph{divisor}.
\end{definition}

\begin{claim}\label{clm:index-returns}
Index returns are the capitalisation-weighted average of constituent returns; the divisor sets only
the level.
\end{claim}

\begin{proof}
Both halves follow from $I_t$ being a \emph{sum divided by a constant}. A constituent's dollar move
is $N_i$ times its price change, so its share of the index move is its share of total
capitalisation:
\[
  r_I = \sum_i w_{it} r_i, \qquad w_{it} = \frac{N_i P_{it}}{M_t}, \qquad \sum_i w_{it} = 1 .
\]
And $D_t$ divides $I_t$ at both dates, so it cancels from the ratio---it can shift the level
anywhere without touching a single return.
\end{proof}

\begin{note}[Two consequences]
\textbf{The top names dominate.} At $w \approx 0.07$ against $w \approx 0.0002$, the largest
constituent carries roughly \textbf{350 times} the index impact of the smallest; this is not an
average of 500 \emph{prices}. \textbf{And the level is arbitrary.} $D_t$ cancels from returns, so
the scale is whatever the base period set it to---for the S\&P 500, 1941--43 $=$ 10.
\textbf{Compare returns, never levels.}
\end{note}

A price level is meaningful only against that same series' own history. Across tickers it carries no
information at all---what each ETF's share price happens to be was fixed by how the fund was
structured at inception. Dividing every series by its own first observation removes that
arbitrariness:

\begin{center}
\begin{tabular}{|l|r|r|r|r|}
\hline
\textbf{Ticker} & \textbf{Raw close, start} & \textbf{Raw close, end} & \textbf{Rebased} & \textbf{Return} \\ \hline
SPY & 325.05 & 740.83 & 227.9 & \textbf{$+127.9\%$} \\ \hline
TLT & 137.10 &  87.35 &  63.7 & \textbf{$-36.3\%$}  \\ \hline
GLD & 143.97 & 368.49 & 255.9 & \textbf{$+155.9\%$} \\ \hline
\end{tabular}
\end{center}

SPY has the highest closing price on every single day of the sample, yet GLD returned more. The end
level alone cannot tell you this; you need the start it is measured against.

\fig{levels-vs-rebased}{The same three ETFs, twice. Left, raw closing prices: SPY runs highest for
the entire sample. Right, each series divided by its own first value and multiplied by 100: GLD now
finishes highest and TLT ends below the line it started on. The ordering of the levels and the
ordering of the returns are different orderings.}{fig:levels-vs-rebased}

\subsection{Why futures rather than ETFs}\label{sec:why-futures}

\begin{tabularx}{\linewidth}{|>{\raggedright\arraybackslash}p{0.20\linewidth}|X|X|}
\hline
\textbf{Reason} & \textbf{Futures} & \textbf{ETFs} \\ \hline
\textbf{Symmetry}
  & Shorting costs nothing extra
  & Shorting needs borrowed shares, a fee, availability \\ \hline
\textbf{Coverage}
  & Deep markets in all four asset classes
  & Poor in commodities and FX \\ \hline
\textbf{Capital efficiency}
  & Margin is roughly 5 to 10 percent of notional; spare cash earns \textbf{collateral yield}
  & Full notional tied up \\ \hline
\end{tabularx}

Symmetry is decisive: a strategy that must go short as freely as long cannot tolerate the
asymmetry, and \secref{sec:short-fragile} shows why the short leg is the fragile one. The cost is
that futures expire, so a continuous series has to be stitched together from
contracts---see \chapref{100 \midot{} The Dataset}{100_dataset}.


\section{Momentum, the Source of the Return}\label{sec:momentum}

\subsection{What momentum is}\label{sec:what-momentum-is}

\begin{definition}[Momentum]\label{def:momentum}
\emph{Momentum} is the empirical tendency for recent relative performance to persist: assets that
outperformed over the past one to twelve months tend to keep outperforming, and recent laggards tend
to keep lagging.
\end{definition}

\begin{note}
The claim is about \emph{returns}, not price levels---it concerns the ordering of performance across
assets and across time. That is why one momentum rule can be applied to markets with nothing else
in common, from Treasury futures to natural gas.
\end{note}

Turning this into a number---a lookback window, a normalization, a test that it carries
information---is \chapref{2 \midot{} Testing a Signal}{02_testing_a_signal}'s subject. Here it is
only the premise.

\subsection{Why momentum might work}\label{sec:why-momentum-works}

Two contested explanations, no proof. Momentum has been durably profitable regardless.

\begin{itemize}
  \item \textbf{Information diffusion.} Unglamorous macro news spreads gradually, creating sustained
        pressure. Possibly weakened since 2008, as computing and data distribution got faster.
  \item \textbf{Selective momentum capture.} Winners keep winning and losers keep losing, because
        large positions take time to build and unwind, investors chase winners, and major events do
        not resolve in a single day.
\end{itemize}

\chapref{2 \midot{} Testing a Signal}{02_testing_a_signal} turns this from a story into something
testable.


\section{Long and Short, the Two Sides of a Position}\label{sec:long-and-short}

\subsection{What the two sides mean}\label{sec:two-sides}

\begin{definition}[Long]\label{def:long}
Buy first, sell later. Profit is the sell price minus the buy price.
\end{definition}

\begin{definition}[Short]\label{def:short}
Borrow and sell first, buy back and return later. Profit is the sell price minus the buy-back price.
\end{definition}

\begin{note}
Running both deploys capital fully---a 150/50 or 200/100 book rather than idle cash---and lets the
strategy express a negative view, which a long-only book cannot.
\end{note}

\subsection{How a short is borrowed}\label{sec:how-a-short-is-borrowed}

A short is not selling out of thin air---the shares are borrowed before they are sold:

\begin{enumerate}
  \item \textbf{Borrow} $N$ shares from a lender---a broker, or a long-term holder.
  \item \textbf{Sell} them now, and receive cash.
  \item \textbf{Buy back} $N$ shares later, from the market.
  \item \textbf{Return} the $N$ shares, closing the position.
\end{enumerate}

\begin{note}
\textbf{You never create shares.} They are fungible, so returning ``$N$ shares of SPY''---not
specific certificates---settles the debt. \textbf{Lenders are paid.} Long-term holders earn a
\textbf{lending fee} on stock they were going to sit on anyway. \textbf{And the broker sits in the
middle}, matching the two sides and holding margin; the market for this is \textbf{securities
lending}.
\end{note}

\pointer{The full mechanics beneath this---who holds the shares mid-loan, who receives the dividend,
hard-to-borrow rates, and how a short squeeze runs away: \S 10 of the
\texttt{market-101-foundations} note in \texttt{market\_knowledge/}.}

\subsection{Why the short leg is the fragile one}\label{sec:short-fragile}

\begin{claim}\label{clm:short-unbounded}
For a static position, a long's loss is bounded by its outlay; a short's is not.
\end{claim}

\begin{proof}
With signed share count $s$, entry price $P_0 > 0$ and exit price $P_T$, the payoff is
\[
  \mathrm{PnL}(P_T) = s \left( P_T - P_0 \right), \qquad P_T \in [0, \infty)
\]
the domain being half-open because a price cannot go negative.

For a \textbf{long}, $s > 0$: $\mathrm{PnL}$ is increasing in $P_T$, so its infimum is at the left
endpoint, $\mathrm{PnL} \geq s(0 - P_0) = -s P_0$---exactly the amount paid. For a \textbf{short},
$s < 0$: $\mathrm{PnL}$ is decreasing in $P_T$, and the domain has no right endpoint, so
$\mathrm{PnL} \to -\infty$ as $P_T \to \infty$.
\end{proof}

The asymmetry is in the \textbf{domain}, not the payoff---both are lines of slope $\left| s \right|$.
The floor exists only because prices are floored at zero.

\fig{payoff-asymmetry}{Profit and loss per dollar of exposure against the terminal price $P_T$. The
long and short payoffs are lines of equal and opposite slope crossing at the entry price $P_0$; the
long's line stops at $-100$ percent where the price reaches zero, while the short's continues
downward without limit as the price rises. Illustrative.}{fig:payoff-asymmetry}

\begin{note}
Hence margin requirements and forced buy-ins---and why a 150/50 book, 200 percent gross, is
\textbf{two times leverage} rather than free money.
\end{note}

\subsection{How a short is represented}\label{sec:how-a-short-is-represented}

A short needs no separate machinery. Carry the position as a \textbf{signed quantity} and let the
sign do the work: a negative share count valued at the market price is a \textbf{negative position
value}, which rises exactly when the price falls.

\begin{note}[Sign convention]
A negative position value is a \textbf{liability}---a debt owed---not a loss already taken.
\end{note}

\begin{note}[Static versus constant-exposure shorts]
The proof above holds the share count fixed at entry. A book that instead maintains a constant
\emph{dollar} exposure buys the short back down as it moves against it, so the unbounded loss is
never realised. That daily rebalancing is an \textbf{implicit risk control}: the proof is right, but
a constant-exposure book is not running the position it describes.
\end{note}

\pointer{How this project implements it, with the accounting derivation and the measured gap between
the two: \texttt{Backtests.md} in \texttt{Backtest\_prototype/}.}


\section{Who Else Is Trading}\label{sec:who-else}

\begin{tabularx}{\linewidth}{|>{\raggedright\arraybackslash}p{0.26\linewidth}|X|>{\raggedright\arraybackslash}p{0.28\linewidth}|}
\hline
\textbf{Participant} & \textbf{Examples} & \textbf{Typical hold} \\ \hline
Index / passive funds & Vanguard, BlackRock & Years to decades \\ \hline
Active managers & Fidelity, PIMCO & Monthly to quarterly \\ \hline
Hedge funds & --- & Intraday to several days \\ \hline
Market makers & Optiver, Citadel Securities, IMC, SIG & Seconds to minutes \\ \hline
Noise traders & Retail, uninformed flow & No consistent horizon \\ \hline
\end{tabularx}


\section*{Background}
\addcontentsline{toc}{section}{Background}

General finance knowledge needed to read this chapter, not part of its argument.

\subsection*{What is the relationship between futures/options and asset classes?}

They answer different questions, so they are not alternatives to one another.

\begin{itemize}
  \item \textbf{Asset class}---\emph{what} underlying market you are exposed to.
  \item \textbf{Instrument}---\emph{how} you hold that exposure: cash, futures, or options.
\end{itemize}

Futures are not an asset class, and options are not a rival one. Every asset class trades in every
instrument:

\begin{tabularx}{\linewidth}{|>{\raggedright\arraybackslash}p{0.20\linewidth}|X|>{\raggedright\arraybackslash}p{0.20\linewidth}|>{\raggedright\arraybackslash}p{0.20\linewidth}|}
\hline
\textbf{Asset class} & \textbf{Typical underlying} & \textbf{Futures} & \textbf{Options} \\ \hline
\textbf{Equities}     & S\&P 500, Nasdaq 100      & ES, NQ                       & options on ES \\ \hline
\textbf{Fixed income} & US Treasuries, Bunds      & ZN (10-year), ZB (30-year)   & options on ZN \\ \hline
\textbf{Commodities}  & Crude oil, gold, corn     & CL, GC, ZC                   & options on CL \\ \hline
\textbf{Currencies (FX)} & EUR/USD, USD/JPY          & 6E, 6J                       & options on 6E \\ \hline
\end{tabularx}

The same exposure comes in three wrappers: equity exposure is bought in cash as SPY, in futures as
ES, or in options as an ES call. So in the opening definition---``trades \textbf{futures}
systematically \textbf{across asset classes}''---the three parts are independent:
\textbf{futures} is the wrapper, \textbf{across asset classes} is the breadth of underlying markets,
and \textbf{systematically} means by rule or model rather than discretion.

\subsection*{Which asset class does a bond belong to?}

\textbf{Fixed income.} Treasuries, corporate, municipal and high-yield bonds all sit there.

\begin{itemize}
  \item ``Fixed'' describes the \textbf{contract}, meaning a scheduled coupon and the return of
        principal---not a fixed price.
  \item Bond prices still move, with interest rates, credit risk and time to maturity. That
        variation is what makes them a trend market at all.
  \item Treasury futures trade at 2-, 5-, 10- and 30-year maturities, so a CTA picks the
        \emph{maturity} it wants exposure to.
\end{itemize}

\subsection*{Why do CTAs use futures rather than options?}

\begin{tabularx}{\linewidth}{|>{\raggedright\arraybackslash}p{0.22\linewidth}|X|X|}
\hline
 & \textbf{Futures} & \textbf{Options} \\ \hline
Exposure
  & \textbf{Linear}---profit and loss move one-for-one with the underlying
  & Non-linear; depends on strike and moneyness \\ \hline
Extra drivers & None & Implied volatility, time decay, strike \\ \hline
What you are betting on & Direction & Direction \textbf{and} volatility \\ \hline
\end{tabularx}

An option position is partly a volatility trade whether you intended it or not. Linear exposure is
what lets one risk framework size 37 positions on the same scale.

\begin{example}\label{ex:four-legs}
A trend follower might simultaneously hold long ES, short ZN, long GC, short 6J---all four are
futures, and all four are different asset classes.
\end{example}


\section*{Appendix \midot{} Notation}
\addcontentsline{toc}{section}{Appendix \midot{} Notation}

Throughout, $t$ is the date and $i$ indexes an index's constituents.

\begin{xltabular}{\linewidth}{|>{\raggedright\arraybackslash}p{0.20\linewidth}|X|>{\raggedright\arraybackslash}p{0.22\linewidth}|}
\hline
\textbf{Symbol} & \textbf{Means} & \textbf{First used} \\ \hline
\endhead
$I_t$
  & the index level on date $t$---a published number, not a tradeable object
  & \secref{sec:index-etf-future} \\ \hline
$P_{it}$, $N_i$
  & constituent $i$'s price on date $t$, and its float-adjusted share count
  & \secref{sec:index-etf-future} \\ \hline
$M_t$, $D_t$
  & the total float capitalisation, and the divisor that turns it into the printed level
  & \secref{sec:index-etf-future} \\ \hline
$w_{it}$
  & constituent $i$'s share of total capitalisation---its weight in the index return
  & \secref{sec:index-etf-future} \\ \hline
$r_I$, $r_i$
  & the index's return and constituent $i$'s return, over the same period
  & \secref{sec:index-etf-future} \\ \hline
$s$
  & the \textbf{signed} share count of a position: positive is long, negative is short
  & \secref{sec:short-fragile} \\ \hline
$P_0$, $P_T$
  & the entry and exit price of that position
  & \secref{sec:short-fragile} \\ \hline
$\mathrm{PnL}(P_T)$
  & the position's profit and loss as a function of where the price ends up
  & \secref{sec:short-fragile} \\ \hline
$N$
  & the number of shares borrowed to open a short
  & \secref{sec:how-a-short-is-borrowed} \\ \hline
\end{xltabular}

\begin{note}[Collisions to watch]
$s$ here is a \textbf{signed share count}, not
\chapref{2}{02_testing_a_signal}'s $s$, which indexes the asset. $w_{it}$ is a
\textbf{constituent's weight inside an index}, fixed by capitalisation and not chosen by anyone;
\chapref{2}{02_testing_a_signal}'s $w_{s,t}$ and
\chapref{5 \midot{} \S 3.3}{05_modelling}'s $w_i$ are \textbf{portfolio} weights, which are decisions.
$r_i$ is a constituent's return, whereas elsewhere in the course $r$ carries an asset and a date
subscript and means the forward return being predicted.
\end{note}
